\section{Distribution Resilience Assessment with Mobile Energy Storage Systems}
\label{sec:distribution-resilience-mess}

\subsection{Problem Form and Runtime Scope}

This chapter documents the resilience-restoration layer for multi-period outage
recovery with mobile energy storage systems.  The mathematical model couples
feeder availability, reconfiguration, served-load restoration, and inter-temporal
MESS movement and discharge decisions over a discrete restoration horizon.

\subsection{Code Map}

The current implementation is concentrated in:
\begin{itemize}
  \item \texttt{src/analysis/resilience\_assessment.cpp} for restoration simulation and resilience KPI aggregation,
  \item \texttt{src/analysis/topology\_analysis.cpp} for feeder-switching support,
  \item distribution and reliability-side result structures used by the GUI and regression tests.
\end{itemize}

\subsection{Pseudocode Map}

The detailed algorithm chapter does not yet include a dedicated MESS-resilience
pseudocode block.  The closest reusable workflow references are the ONR and
distribution-PF algorithms (A13 and A14) in Section~\ref{sec:algorithms}.  This
chapter therefore remains the main code-aligned statement of the restoration
model itself.

\subsection{Current Status}

The C++ port now has two layers. The default runtime path remains the practical
heuristic restoration simulator used by the GUI and validation demo, while a
new strict multi-period MIP/LinDistFlow restoration skeleton has been added as
the replacement path for unexpected-fault studies. The current C++ scope is
therefore broader than before, but still narrower than the original full
stochastic two-stage decomposition design.

This section documents the resilience-restoration model used for
multi-period distribution-service recovery with mobile energy storage systems
(MESS). The formulation is aligned with the original MATLAB package in
\texttt{Multi-stage Restoration/} and with the current C++ implementation in
\texttt{resilience\_assessment.cpp}.

The original MATLAB workflow is a two-stage stochastic mixed-integer recourse
model with Benders-style decomposition. The current C++ port implements the
core operational layer needed for practical resilience studies in this codebase:
\begin{enumerate}
  \item time-indexed feeder outage and repair simulation,
  \item radial reconfiguration through healthy tie-line closure,
  \item island-level load restoration prioritization,
  \item inter-temporal MESS routing and discharge scheduling,
  \item resilience KPI aggregation for GUI and regression tests.
\end{enumerate}

In addition, the latest C++ implementation now contains a strict restoration
builder that formulates a single multi-period MIP with binary branch-status,
energized-bus, and root-selection variables; active-power LinDistFlow voltage
envelopes; fixed-storage inter-temporal energy balance; and time-space MESS
movement arcs. This strict model is currently a first skeleton, intended to
replace the heuristic unexpected-fault restoration path incrementally.

\subsection{Sets and Indices}

\begin{align}
\mathcal{N} &= \text{distribution buses}, \\
\mathcal{E} &= \text{AC feeder branches}, \\
\mathcal{E}^{\mathrm{tie}} &\subseteq \mathcal{E} && \text{normally open tie branches}, \\
\mathcal{G} &= \text{dispatchable local generators}, \\
\mathcal{R} &= \text{static and renewable distributed resources}, \\
\mathcal{S} &= \text{fixed stationary storage units}, \\
\mathcal{M} &= \text{mobile energy storage systems (MESS)}, \\
\mathcal{L} &= \text{loads}, \\
\mathcal{T} &= \{0,1,\dots,T-1\} && \text{discrete restoration periods}, \\
\mathcal{F} &= \text{faulted feeder branches}. 
\end{align}

Each load $\ell\in\mathcal{L}$ is attached to a bus $n(\ell)\in\mathcal{N}$,
with active demand $P^{d}_{\ell,t}$, priority weight $w_{\ell}$, and customer
count $N^{\mathrm{cust}}_{\ell}$.

\subsection{State Variables}

For each time period $t\in\mathcal{T}$:
\begin{align}
\alpha_{e,t} &\in \{0,1\} && \text{branch energization state}, \\
P^{\mathrm{serv}}_{\ell,t} &\ge 0 && \text{served load}, \\
P^{\mathrm{shed}}_{\ell,t} &\ge 0 && \text{unsupplied load}, \\
E^{\mathrm{mess}}_{m,t} &\ge 0 && \text{MESS stored energy}, \\
P^{\mathrm{mess}}_{m,t} &\ge 0 && \text{MESS discharge power}, \\
\chi^{\mathrm{stay}}_{m,n,t} &\in \{0,1\} && \text{MESS located at bus }n, \\
\chi^{\mathrm{move}}_{m,a,t} &\in \{0,1\} && \text{MESS traverses transport arc }a.
\end{align}

The C++ implementation reports these decisions in aggregated form through the
per-step result objects:
\begin{align}
P^{\mathrm{demand}}_t,\;
P^{\mathrm{served}}_t,\;
P^{\mathrm{shed}}_t,\;
R_t,\;
\mathcal{E}^{\mathrm{open}}_t,\;
\mathcal{M}^{\mathrm{dispatch}}_t.
\end{align}

\subsection{Fault and Repair Process}

A faulted branch $f\in\mathcal{F}$ is unavailable during its outage window:
\begin{align}
\alpha_{f,t} = 0,
\qquad
\forall t\in\left[t^{\mathrm{out}}_f,\; t^{\mathrm{out}}_f + \tau^{\mathrm{rep}}_f\right).
\end{align}

Outside this interval, the branch may either return to its original state or be
closed by reconfiguration if it is a healthy tie branch. In the deterministic
C++ implementation, repair times are explicit inputs and branch failures are
simulated chronologically over the horizon.

\subsection{Load Restoration and Prioritization}

For every load $\ell$ and time $t$:
\begin{align}
P^{\mathrm{serv}}_{\ell,t} + P^{\mathrm{shed}}_{\ell,t} = P^{d}_{\ell,t},
\qquad
0 \le P^{\mathrm{serv}}_{\ell,t} \le P^{d}_{\ell,t}.
\end{align}

Within each island, available supply is allocated to higher-priority loads
first. If $\pi(1),\pi(2),\dots$ denotes the ordering of island loads by
descending $w_{\ell}$, then the restoration rule is
\begin{align}
P^{\mathrm{serv}}_{\pi(k),t}
= \min\left(P^{d}_{\pi(k),t},
\; P^{\mathrm{avail}}_{t} - \sum_{j=1}^{k-1} P^{\mathrm{serv}}_{\pi(j),t}\right)_{+}.
\end{align}

The weighted unserved load is
\begin{align}
P^{\mathrm{wshed}}_t = \sum_{\ell\in\mathcal{L}} w_{\ell} P^{\mathrm{shed}}_{\ell,t}.
\end{align}

\subsection{Island Supply Model}

For each electrical island $\mathcal{I}\subseteq\mathcal{N}$ at time $t$, the
available active supply is approximated as
\begin{align}
P^{\mathrm{avail}}_{\mathcal{I},t} =
P^{\mathrm{grid}}_{\mathcal{I},t}
+ \sum_{g\in\mathcal{G}(\mathcal{I})} P^{\max}_{g}
+ \sum_{r\in\mathcal{R}(\mathcal{I})} P_{r,t}
+ \sum_{s\in\mathcal{S}(\mathcal{I})} P^{\mathrm{dis}}_{s,t}
+ \sum_{m\in\mathcal{M}(\mathcal{I})} P^{\mathrm{mess}}_{m,t}.
\end{align}

If an island contains a slack or external-grid source, then
$P^{\mathrm{grid}}_{\mathcal{I},t}$ is treated as nonbinding and the island can
serve its full demand. Otherwise restoration is limited by local DER and
storage capability.

\subsection{MESS Energy Dynamics}

For each mobile unit $m\in\mathcal{M}$:
\begin{align}
E^{\min}_{m} \le E^{\mathrm{mess}}_{m,t} \le E^{\max}_{m},
\end{align}
and during a discharge period,
\begin{align}
E^{\mathrm{mess}}_{m,t+1}
= E^{\mathrm{mess}}_{m,t}
- \frac{\Delta t}{\eta^{\mathrm{dis}}_{m}} P^{\mathrm{mess}}_{m,t}.
\end{align}

Travel consumes energy according to the traversed road distance:
\begin{align}
E^{\mathrm{mess}}_{m,t^{+}}
= E^{\mathrm{mess}}_{m,t^{-}} - \rho^{\mathrm{travel}}_m d_{m,t},
\end{align}
where $\rho^{\mathrm{travel}}_m$ is the transport energy cost in MWh/km and
$d_{m,t}$ is the routing distance during the dispatch decision.

The discharge capability is limited by both power and stored energy:
\begin{align}
0 \le P^{\mathrm{mess}}_{m,t}
\le \min\left(P^{\max}_{m},
\frac{\eta^{\mathrm{dis}}_{m}}{\Delta t}
\left(E^{\mathrm{mess}}_{m,t} - E^{\min}_{m}\right)\right).
\end{align}

\subsection{Transportation-Network Coupling}

The original MATLAB model defines a time-space transportation network with
binary path-flow variables. The corresponding canonical constraints are:
\begin{align}
\sum_{a\in\delta^{+}(v)} \chi^{\mathrm{move}}_{m,a,t}
- \sum_{a\in\delta^{-}(v)} \chi^{\mathrm{move}}_{m,a,t}
= b_{m,v,t},
\end{align}
where $b_{m,v,t}$ is the source/sink balance induced by the current MESS
location and the desired destination.

In the current C++ implementation, this is approximated by a shortest-path
assignment over a transport graph. If no explicit road network is provided, the
AC feeder graph is used as a transportation proxy with branch length as edge
weight. This keeps the data contract lightweight while preserving the essential
spatial coupling between restoration need and MESS mobility.

\subsection{Radial Reconfiguration}

Healthy open ties may be closed to reconnect unsupplied islands, subject to a
radial forest policy. In the simplified implementation, branch closures are
accepted only when they connect two different components. This enforces
cycle-free restoration and mimics the spanning-tree intent of the MATLAB master
problem:
\begin{align}
\alpha_{e,t} = 1
\implies e \text{ joins two distinct connected components at time } t.
\end{align}

This heuristic is consistent with single-commodity radial restoration logic and
remains the default GUI path. In parallel, the code base now also contains a
strict mixed-integer LinDistFlow restoration skeleton, so future unexpected-
fault studies no longer need to remain confined to heuristic topology updates.

\subsection{Objective Function}

The stochastic MATLAB model minimizes expected restoration cost across
scenarios. Its operational core is represented by
\begin{align}
\min \sum_{t\in\mathcal{T}} \Bigg[
C^{\mathrm{shed}} \sum_{\ell\in\mathcal{L}} P^{\mathrm{shed}}_{\ell,t}
+ C^{\mathrm{sw}} \sum_{e\in\mathcal{E}^{\mathrm{tie}}}
\left|\alpha_{e,t} - \alpha_{e,t-1}\right|
+ C^{\mathrm{travel}} \sum_{m\in\mathcal{M}} d_{m,t}
\Bigg].
\end{align}

The implemented C++ routine uses the same operational priorities, but solves
them sequentially by:
\begin{enumerate}
  \item enforcing feeder outages and completed repairs,
  \item greedily reconnecting unsourced islands through healthy ties,
  \item allocating local generation and storage to maximize served load,
  \item dispatching available MESS to the most critical unsupplied islands.
\end{enumerate}

\subsection{Resilience Metrics}

The primary resilience outputs are:
\begin{align}
E^{\mathrm{demand}} &= \sum_{t\in\mathcal{T}} P^{\mathrm{demand}}_t \Delta t, \\
E^{\mathrm{served}} &= \sum_{t\in\mathcal{T}} P^{\mathrm{served}}_t \Delta t, \\
E^{\mathrm{uns}} &= \sum_{t\in\mathcal{T}} P^{\mathrm{shed}}_t \Delta t, \\
R_t &= \frac{P^{\mathrm{served}}_t}{P^{\mathrm{demand}}_t}, \\
\mathrm{RI} &= \frac{E^{\mathrm{served}}}{E^{\mathrm{demand}}}, \\
E^{\mathrm{MESS}} &= \sum_{t\in\mathcal{T}} \sum_{m\in\mathcal{M}} P^{\mathrm{mess}}_{m,t} \Delta t.
\end{align}

The GUI reports these quantities as:
\begin{itemize}
  \item total served energy,
  \item total shed energy,
  \item resilience index,
  \item average and final restoration ratio,
  \item peak shed power,
  \item MESS-delivered energy,
  \item MESS travel distance,
  \item switching actions and repaired-fault count.
\end{itemize}

\subsection{Implementation Notes}

The present C++ port is intentionally conservative in scope:
\begin{enumerate}
  \item it runs directly on the canonical \texttt{HybridPowerSystem} data model,
  \item it supports arbitrary user-specified feeder faults and repair windows,
  \item it uses explicit mobile-storage state tracking over time,
  \item it returns lightweight JSON-friendly outputs for the GUI,
  \item it remains extensible toward the full stochastic Benders architecture in
        the original MATLAB package.
\end{enumerate}

This design is appropriate for interactive resilience screening and regression
validation, while preserving a clean upgrade path toward a full mixed-integer
recourse implementation.

\subsection{Validation Snapshot}

The new regression test \texttt{tests/test\_resilience\_assessment.cpp} uses the
\texttt{case33mg\_acdc} feeder with a single automatically selected feeder
fault, a six-hour repair window, and an eight-hour restoration horizon. To
isolate the marginal contribution of MESS, stationary DER is removed from the
test scenario and one mobile unit is placed on the tail bus.

The measured results are:
\begin{align}
E^{\mathrm{demand}} &= 29.72\ \text{MWh}, \\
E^{\mathrm{shed,base}} &= 0.48\ \text{MWh}, \\
E^{\mathrm{shed,MESS}} &= 0.00\ \text{MWh}, \\
\mathrm{RI}^{\mathrm{base}} &= 0.983849, \\
\mathrm{RI}^{\mathrm{MESS}} &= 1.000000, \\
E^{\mathrm{MESS}} &= 0.48\ \text{MWh}.
\end{align}

In other words, the mobile storage unit fully eliminates the residual outage
energy left by the feeder fault in this validation case. This is the expected
behavior for a tail-island restoration event and provides a useful sanity check
for both the operational logic and the GUI integration.
